\chapter{Préparation des Données et Modélisation}

\begin{objectivebox}
Ce chapitre combine les phases CRISP-DM de préparation des données et de modélisation, détaillant la conception de l'architecture technique, le développement des fonctionnalités et l'implémentation de la solution Housy Tunisia.
\end{objectivebox}

\section{Préparation des Données (Data Preparation)}

\subsection{Nettoyage et Transformation}

\subsubsection{Stratégies de Nettoyage}

Les données collectées ont nécessité plusieurs opérations de nettoyage :

\begin{itemize}
    \item \textbf{Suppression des doublons :} 12 projets identiques supprimés
    \item \textbf{Correction des valeurs aberrantes :} Validation manuelle de 45 cas
    \item \textbf{Standardisation des formats :} Harmonisation des unités (TND, m², etc.)
    \item \textbf{Validation géographique :} Correction de 28 coordonnées GPS erronées
    \item \textbf{Imputation des valeurs manquantes :} Algorithmes KNN pour 156 cas
\end{itemize}

\textbf{[CAPTURE À AJOUTER : Interface de nettoyage des données avec statistiques]}

\subsubsection{Transformation des Données}

\begin{itemize}
    \item \textbf{Normalisation :} Standardisation des échelles pour les variables numériques
    \item \textbf{Encodage catégoriel :} One-hot encoding pour les types de construction
    \item \textbf{Création de variables dérivées :} Prix/m², ratios financiers
    \item \textbf{Agrégation temporelle :} Moyennes mobiles sur les prix des matériaux
\end{itemize}

\subsection{Structuration en Base de Données}

\subsubsection{Conception du Schéma}

Le schéma PostgreSQL a été optimisé pour :

\begin{itemize}
    \item \textbf{Performance :} Index sur les champs fréquemment requis
    \item \textbf{Intégrité :} Contraintes de clés étrangères
    \item \textbf{Extensibilité :} Structure modulaire pour ajouts futurs
    \item \textbf{Historisation :} Suivi des modifications temporelles
\end{itemize}

\begin{lstlisting}[language=SQL, caption=Exemple de table optimisée]
CREATE TABLE projects (
    id SERIAL PRIMARY KEY,
    name VARCHAR(255) NOT NULL,
    project_type project_type_enum NOT NULL,
    total_surface DECIMAL(10,2),
    estimated_cost DECIMAL(15,2),
    location_id INTEGER REFERENCES locations(id),
    created_at TIMESTAMP DEFAULT CURRENT_TIMESTAMP,
    updated_at TIMESTAMP DEFAULT CURRENT_TIMESTAMP
);

CREATE INDEX idx_projects_type ON projects(project_type);
CREATE INDEX idx_projects_location ON projects(location_id);
\end{lstlisting}

\textbf{[CAPTURE À AJOUTER : Diagramme ERD complet de la base de données]}

\section{Architecture Technique (Modeling)}

\subsection{Stack Technologique}

\subsubsection{Frontend : React.js + TypeScript}

L'interface utilisateur a été développée avec :

\begin{itemize}
    \item \textbf{React 18 :} Framework principal avec hooks modernes
    \item \textbf{TypeScript :} Typage statique pour la robustesse
    \item \textbf{Vite :} Outil de build rapide et moderne
    \item \textbf{Tailwind CSS :} Framework CSS utility-first
    \item \textbf{React Router :} Gestion du routing côté client
\end{itemize}

\begin{lstlisting}[language=TypeScript, caption=Composant React typé]
interface ProjectCardProps {
    project: {
        id: number;
        name: string;
        type: ProjectType;
        estimatedCost: number;
        surface: number;
    };
    onEdit: (id: number) => void;
}

const ProjectCard: React.FC<ProjectCardProps> = ({ project, onEdit }) => {
    return (
        <div className="bg-white rounded-lg shadow-md p-6">
            <h3 className="text-xl font-semibold mb-2">{project.name}</h3>
            <p className="text-gray-600">Type: {project.type}</p>
            <p className="text-gray-600">Surface: {project.surface} m²</p>
            <p className="text-green-600 font-bold">
                Coût estimé: {project.estimatedCost.toLocaleString()} TND
            </p>
            <button 
                onClick={() => onEdit(project.id)}
                className="mt-4 bg-blue-500 text-white px-4 py-2 rounded hover:bg-blue-600"
            >
                Modifier
            </button>
        </div>
    );
};
\end{lstlisting}

\textbf{[CAPTURE À AJOUTER : Interface principale de l'application]}

\subsubsection{Backend : Node.js + Express.js}

Le serveur backend utilise :

\begin{itemize}
    \item \textbf{Node.js 18+ :} Runtime JavaScript haute performance
    \item \textbf{Express.js :} Framework web minimaliste et flexible
    \item \textbf{TypeScript :} Cohérence avec le frontend
    \item \textbf{PostgreSQL :} Base de données relationnelle robuste
    \item \textbf{Redis :} Cache et gestion des sessions
\end{itemize}

\begin{lstlisting}[language=TypeScript, caption=API REST pour la gestion des projets]
import express from 'express';
import { ProjectService } from '../services/ProjectService';

const router = express.Router();
const projectService = new ProjectService();

// GET /api/projects - Liste des projets
router.get('/', async (req, res) => {
    try {
        const { page = 1, limit = 10, type } = req.query;
        const projects = await projectService.getProjects({
            page: Number(page),
            limit: Number(limit),
            type: type as string
        });
        res.json(projects);
    } catch (error) {
        res.status(500).json({ error: 'Erreur serveur' });
    }
});

// POST /api/projects - Création de projet
router.post('/', async (req, res) => {
    try {
        const projectData = req.body;
        const newProject = await projectService.createProject(projectData);
        res.status(201).json(newProject);
    } catch (error) {
        res.status(400).json({ error: 'Données invalides' });
    }
});

export { router as projectRoutes };
\end{lstlisting}

\textbf{[CAPTURE À AJOUTER : Architecture technique globale de l'application]}

\subsection{Fonctionnalités Développées}

\subsubsection{Gestion des Projets}

Les fonctionnalités de gestion de projets incluent :

\begin{itemize}
    \item \textbf{CRUD complet :} Création, lecture, mise à jour, suppression
    \item \textbf{Recherche avancée :} Filtres par type, statut, budget
    \item \textbf{Planification :} Diagrammes de Gantt interactifs
    \item \textbf{Suivi d'avancement :} Pourcentage de completion en temps réel
    \item \textbf{Gestion des ressources :} Allocation équipes et matériaux
\end{itemize}

\textbf{[CAPTURE À AJOUTER : Interface de gestion des projets avec tableau de bord]}

\subsubsection{Système d'Estimation}

Le moteur d'estimation intègre :

\begin{itemize}
    \item \textbf{Calculs automatisés :} Métrés basés sur les dimensions
    \item \textbf{Base de prix actualisée :} Tarifs matériaux et main-d'œuvre TND
    \item \textbf{Variations régionales :} Adaptation aux spécificités locales
    \item \textbf{Niveaux de finition :} Standard, haut de gamme, luxe
    \item \textbf{Export multi-formats :} PDF, Excel, Word
\end{itemize}

\begin{lstlisting}[language=TypeScript, caption=Service d'estimation]
class EstimationService {
    async calculateProjectCost(projectData: ProjectInput): Promise<Estimation> {
        const { surface, type, location, finishLevel } = projectData;
        
        // Récupération des prix matériaux
        const materialCosts = await this.getMaterialCosts(location);
        
        // Calcul main-d'œuvre
        const laborCosts = await this.getLaborCosts(location, type);
        
        // Application des coefficients de finition
        const finishMultiplier = this.getFinishMultiplier(finishLevel);
        
        // Calcul total
        const baseCost = (materialCosts + laborCosts) * surface;
        const totalCost = baseCost * finishMultiplier;
        
        return {
            totalCost,
            breakdown: {
                materials: materialCosts * surface,
                labor: laborCosts * surface,
                finish: (totalCost - baseCost)
            },
            currency: 'TND',
            createdAt: new Date()
        };
    }
}
\end{lstlisting}

\textbf{[CAPTURE À AJOUTER : Interface d'estimation avec détail des coûts]}

\subsubsection{Gestion des Matériaux}

Le catalogue de matériaux comprend :

\begin{itemize}
    \item \textbf{800+ références :} Matériaux adaptés au marché tunisien
    \item \textbf{Fiches techniques :} Spécifications détaillées
    \item \textbf{Historique des prix :} Évolution temporelle
    \item \textbf{Fournisseurs :} Base de données des fournisseurs locaux
    \item \textbf{Gestion des stocks :} Suivi des inventaires
\end{itemize}

\textbf{[CAPTURE À AJOUTER : Interface de gestion des matériaux et fournisseurs]}

\section{Optimisation des Performances}

\subsection{Cache Redis}

Redis est utilisé pour optimiser les performances :

\begin{itemize}
    \item \textbf{Sessions utilisateurs :} Stockage rapide des sessions
    \item \textbf{Cache de requêtes :} Mise en cache des calculs fréquents
    \item \textbf{Cache de prix :} Stockage temporaire des tarifs
    \item \textbf{Rate limiting :} Protection contre les abus d'API
\end{itemize}

\begin{lstlisting}[language=TypeScript, caption=Implémentation du cache Redis]
import Redis from 'ioredis';

class CacheService {
    private redis = new Redis(process.env.REDIS_URL);
    
    async getEstimation(key: string): Promise<Estimation | null> {
        const cached = await this.redis.get(`estimation:${key}`);
        return cached ? JSON.parse(cached) : null;
    }
    
    async setEstimation(key: string, estimation: Estimation): Promise<void> {
        await this.redis.setex(
            `estimation:${key}`, 
            3600, // 1 heure
            JSON.stringify(estimation)
        );
    }
    
    async clearEstimationCache(pattern: string): Promise<void> {
        const keys = await this.redis.keys(`estimation:${pattern}*`);
        if (keys.length > 0) {
            await this.redis.del(...keys);
        }
    }
}
\end{lstlisting}

\subsection{Optimisation Base de Données}

\begin{itemize}
    \item \textbf{Index optimisés :} Index composites pour requêtes fréquentes
    \item \textbf{Requêtes optimisées :} Utilisation de vues matérialisées
    \item \textbf{Partitionnement :} Tables partitionnées par date
    \item \textbf{Connection pooling :} Gestion efficace des connexions
\end{itemize}

\textbf{[CAPTURE À AJOUTER : Dashboard de monitoring des performances]}

\section{Sécurité et Authentification}

\subsection{Système d'Authentification}

\begin{itemize}
    \item \textbf{JWT Tokens :} Authentification stateless
    \item \textbf{Hachage bcrypt :} Stockage sécurisé des mots de passe
    \item \textbf{Rôles utilisateurs :} Admin, Manager, Utilisateur
    \item \textbf{Gestion des permissions :} Contrôle d'accès granulaire
\end{itemize}

\begin{lstlisting}[language=TypeScript, caption=Middleware d'authentification]
import jwt from 'jsonwebtoken';

interface AuthMiddleware {
    (req: AuthRequest, res: Response, next: NextFunction): void;
}

export const authenticate: AuthMiddleware = (req, res, next) => {
    const token = req.header('Authorization')?.replace('Bearer ', '');
    
    if (!token) {
        return res.status(401).json({ error: 'Token manquant' });
    }
    
    try {
        const decoded = jwt.verify(token, process.env.JWT_SECRET!) as JWTPayload;
        req.user = decoded;
        next();
    } catch (error) {
        res.status(401).json({ error: 'Token invalide' });
    }
};
\end{lstlisting}

\subsection{Sécurisation des Données}

\begin{itemize}
    \item \textbf{Chiffrement HTTPS :} Communication sécurisée
    \item \textbf{Validation des entrées :} Sanitisation des données
    \item \textbf{Protection CSRF :} Tokens anti-forgery
    \item \textbf{Headers de sécurité :} Configuration Helmet.js
\end{itemize}

\textbf{[CAPTURE À AJOUTER : Interface de gestion des utilisateurs et permissions]}

\section{Types de Construction Tunisiens}

\subsection{Intégration des Spécificités Locales}

L'application intègre les types de construction spécifiques à la Tunisie :

\begin{table}[h]
\centering
\begin{tabular}{|l|l|l|l|}
\hline
\textbf{Type} & \textbf{Caractéristiques} & \textbf{Coût/m² Moyen} & \textbf{Durée Type} \\
\hline
Villa Traditionnelle & Architecture arabo-andalouse & 800-1200 TND & 10-14 mois \\
\hline
Villa Moderne & Design contemporain & 900-1400 TND & 8-12 mois \\
\hline
Immeuble R+4 & Standard urbain tunisien & 600-900 TND & 12-18 mois \\
\hline
Logement Social & Programme gouvernemental & 400-600 TND & 6-10 mois \\
\hline
\end{tabular}
\caption{Types de construction tunisiens intégrés}
\end{table}

\textbf{[CAPTURE À AJOUTER : Galerie des types de construction avec estimations]}

\begin{resultbox}
Cette phase de modélisation a abouti au développement d'une application web complète et fonctionnelle, intégrant toutes les spécificités du marché tunisien de la construction. L'architecture technique moderne garantit performance, sécurité et évolutivité.
\end{resultbox}
